
\subsection{Uncertainty Quantification}
\label{app:inference}

We describe how we construct confidence intervals for our estimators.
We linearize the joint estimating equations, adapting the SNMM inference
construction of \citet[Theorem~8]{lewis2021dynamic}.
Each completed trajectory, including its terminal outcome, is treated as
one independent and identically distributed (IID) observation.
Because multiple decisions within a trajectory share the same terminal
outcome, we first sum their scores within each trajectory, then estimate
the covariance across trajectories.


\paragraph{SNMM estimates.}
For a changed continuation $\pi$, let
$m_i(\alpha,\{\gamma_s\};\pi)$ be trajectory $i$'s contribution
to~\eqref{eq:joint-snmm-moments}. We order its coordinates with
the shared coefficients first, followed by the phase-specific
coefficients for $s=0,\ldots,P-1$. For the shared coefficients,
we sum $Z_{it,\mathcal S}\varepsilon_{it}$ over all decisions
in trajectory $i$. For the coefficients specific to phase $s$,
we sum $Z_{it,\mathcal U}\varepsilon_{it}$ only over decisions
in that phase. Here $\varepsilon_{it}$ is the residual
in~\eqref{eq:joint-snmm-residual}.

For randomized continuation, the observed trajectories already
follow $g$ after the current decision. The conditional contrast
in the observed terminal outcome therefore identifies the blip
under $g$ by~\eqref{eq:identification}. We obtain the estimating
equations by differentiating the least-squares objective
in~\eqref{eq:local-snmm}. We form $m_i(\alpha,\{\gamma_s\};g)$ in exactly the same way as in the previous paragraph, with the only difference being the form of the residual $\varepsilon_{it}$, which is of the form:
\[
\varepsilon_{it} = Y_i-\widehat q_{it}-Z_{it}^{\top}\widehat\beta_{s_{it}}(g),
\]
i.e. does not contain any future treatment correction term, since the observed trajectories already follow $g$ after the current decision.

Having definedd the form of the estimating equations $m_i(\alpha,\{\gamma_s\};g)$ in either case, we can now proceed with the inference. For each SNMM coefficient estimator, we hold the cross-fitted nuisance predictions
fixed and differentiate the summed score with respect to the
coefficients. Define
\[
 \widehat J(\pi)
   =-\sum_{i=1}^n
     \left.
       \frac{\partial m_i(\alpha,\{\gamma_s\};\pi)}
       {\partial(\alpha^\top,\gamma_0^\top,\ldots,\gamma_{P-1}^\top)}
     \right|_{\alpha=\widehat\alpha(\pi),\,\gamma_s=\widehat\gamma_s(\pi)}.
\]
The components of
$\widehat J(\pi)^{-1}m_i(\widehat\alpha(\pi),\{\widehat\gamma_s(\pi)\};\pi)$
give trajectory $i$'s contributions (a.k.a. influence functions) to the linearization of
the fitted coefficients. Write $\dot\alpha_i(\pi)$ for the
components corresponding to the shared coefficients and
$\dot\gamma_{i,s}(\pi)$ for those corresponding to phase $s$.
We assemble the contribution for each phase as
\[
 \dot\beta_{i,s,\mathcal S}(\pi)=\dot\alpha_i(\pi),\qquad
 \dot\beta_{i,s,\mathcal U}(\pi)=\dot\gamma_{i,s}(\pi).
\]
For inference on the causal targets, we assume that the effect
model and its coefficient-sharing restrictions are correctly
specified. The population Jacobian must be nonsingular on the
retained feature directions, and the cross-fitted trajectory
scores must satisfy the regularity conditions for the
linearization.




For a group $G$, define
\[
 a_{i,s}(G)
   =\sum_{t\in\mathcal T_i:s_{it}=s}I_{it}(G)f(H_{it}),
 \qquad
 a_s(G)=\sum_i a_{i,s}(G).
\]
Write $\widehat\tau_{\mathrm{SNMM}}(G;\pi)$ for the SNMM group-effect
estimator under continuation policy $\pi$.
In particular, this corresponds to $\widehat\tau_{\mathrm{SNMM}}(G;g)$
for randomized continuation and $\widehat\tau(G;0)$ for no-treatment
continuation, where $\pi=\pi_0$.
The trajectory contribution is captured by the influence function
\begin{equation}
 u_i\{\widehat\tau_{\mathrm{SNMM}}(G;\pi)\}
   =\frac{
      \sum_s a_{i,s}(G)^\top\widehat\beta_s(\pi)
      -N_i(G)\widehat\tau_{\mathrm{SNMM}}(G;\pi)
      +\sum_s a_s(G)^\top\dot\beta_{i,s}(\pi)
     }{N(G)}.
 \label{eq:blip-influence}
\end{equation}

For policy values, define
\[
 d_{i,s}(\pi)
   =\sum_{t\in\mathcal T_i:s_{it}=s}
      \{\pi(H_{it})-A_{it}\}f(H_{it}),
 \qquad
 d_s(\pi)=\sum_i d_{i,s}(\pi).
\]
The value and gain contributions are captured by the influence functions
\begin{equation}
\begin{aligned}
 u_i\{\widehat V(\pi)\}
   &=\frac{
       Y_i+\sum_s d_{i,s}(\pi)^\top\widehat\beta_s(\pi)
       -\widehat V(\pi)
       +\sum_s d_s(\pi)^\top\dot\beta_{i,s}(\pi)
      }{n},\\
 u_i\{\widehat\Delta(G)\}
   &=u_i\{\widehat V(\pi_G)\}
     -u_i\{\widehat V(\pi_0)\}.
\end{aligned}
\label{eq:value-influence}
\end{equation}
Policy gains average over all completed trajectories, including those
without matching opportunities.

\paragraph{AIPW estimates.}
For the blip with randomized continuation, the trajectory contribution (influence) is
\[
 u_i\{\widehat\tau_{\mathrm{AIPW}}(G;g)\}
   =\frac{1}{N(G)}
      \sum_{t\in\mathcal T_i} I_{it}(G)
      \bigl\{\widehat D_{it}
             -\widehat\tau_{\mathrm{AIPW}}(G;g)\bigr\}.
\]

\paragraph{Reported intervals.}
Let $u_i(\widehat\eta)$ denote the trajectory contribution defined above,
already divided by $N(G)$ for group effects or by $n$ for policy values and
gains.
We report pointwise 95\% intervals in our figures:
\[
 \widehat\eta
 \;\pm\;
 t_{0.975,n-1}
 \left\{
   \frac{n}{n-1}\sum_{i=1}^n u_i(\widehat\eta)^2
 \right\}^{1/2}.
\]
